\documentclass[10pt,a4paper,ragged2e,withhyper]{style} \geometry{left=1cm,right=1cm,top=0.4cm,bottom=0.4cm,columnsep=1.2cm} \usepackage{paracol} \usepackage{xcolor} \definecolor{MainGreen}{HTML}{1f2c67} \definecolor{SlateGrey}{HTML}{2E2E2E} \definecolor{LightGrey}{HTML}{666666}  \colorlet{heading}{MainGreen} \colorlet{headingrule}{MainGreen} \colorlet{accent}{MainGreen} \colorlet{emphasis}{SlateGrey} \colorlet{body}{LightGrey} \renewcommand{\cvItemMarker}{{\small\textbullet}} \renewcommand{\cvRatingMarker}{\faCircle} \renewcommand{\cvsectionfont}{\large\bfseries} \renewcommand{\cvsubsectionfont}{\normalsize\bfseries}  \renewcommand{\namefont}{\Huge\rmfamily\bfseries}
% --- Personnalisations de mise en forme propres à ce document (n'affectent
% pas style.cls ni les autres CV) : ligne décorative fine sous les titres de
% section, dates en italique, espacements resserrés pour un rendu compact et
% professionnel, calibrés pour remplir une page A4 en une seule colonne.
\renewcommand{\cvsection}[2][]{%
  \par\vspace{2pt}%
  \ifstrequal{#1}{}{}{\marginpar{\vspace*{\dimexpr1pt-\baselineskip}\raggedright\input{#1}}}%
  {\color{heading}\cvsectionfont \MakeUppercase{#2}}\par%
  \vspace{0.5pt}
  {\color{headingrule}\hdashrule{\linewidth}{0.5pt}{2pt}\par}\vspace{0.5pt}
}
\renewcommand{\cvevent}[4]{%
  \normalsize\par\noindent
  {\bfseries\color{emphasis}#1\ifstrequal{#2}{}{}{ -- #2}}%
  \ifstrequal{#3}{}{}{%
    \hfill{\color{accent}\bfseries\itshape\cvDateMarker~%
      \BeginAccSupp{method=pdfstringdef,ActualText={\datename: }}#3\EndAccSupp{}}%
  }\par
  \ifstrequal{#4}{}{}{%
    {\small\itshape\color{body}\cvLocationMarker~%
      \BeginAccSupp{method=pdfstringdef,ActualText={\locationname: }}#4\EndAccSupp{}\par}}%
  \vspace{0.5pt}%
  \footnotesize
}
\renewcommand{\divider}{\par\noindent\textcolor{body!30}{\hdashrule{\linewidth}{0.5pt}{0.4ex}}\vspace{0.5pt}}
\setlist{leftmargin=*,labelsep=0.5em,nosep,itemsep=0.02\baselineskip,after=\vspace{0.02\baselineskip}}
\begin{document} \name{El Moustpha Cheikh JIDDOU} \tagline{\textcolor{MainGreen}{Analyste Macroéconomique -- Modélisation DSGE, Politique Monétaire et Comptabilité Nationale}} \personalinfo{ \email{elmoustapha.cheikh.jiddou@gmail.com} \phone{+22248647376} \location{Nouakchott, Mauritanie} \printinfo{\faLinkedin}{El Moustpha Cheikh Jiddou}[https://www.linkedin.com/in/el-moustapha-cheikh-jidou/] \printinfo{\faGithub}{MoustaphaCheikhJidou}[https://github.com/MoustaphaCheikhJidou] \printinfo{\faGlobe}{Portfolio}[https://moustaphacheikhjidou.github.io/MyPortfolio/] } \makecvheader \columnratio{0.6} \begin{paracol}{2}

\cvsection{Expériences Professionnelles}
\cvevent{Stage de fin d'études -- Analyse macroéconomique et modélisation DSGE}{Banque Centrale de Mauritanie (BCM)}{Février 2026 -- Mai 2026}{Nouakchott, Mauritanie}
\begin{itemize}
  \item Réalisation d'un mémoire de master en Économétrie et Statistique Appliquée comparant les régimes monétaires en Mauritanie à l'aide d'un modèle DSGE de petite économie ouverte.
  \item Analyse des mécanismes de transmission de la politique monétaire, de l'inflation, du taux de change et des dynamiques macroéconomiques dans le contexte mauritanien.
  \item Participation aux activités de la Direction Générale des Études et de la Stabilité Monétaire : analyse conjoncturelle, prévisions macroéconomiques et statistiques monétaires.
  \item \textbf{Environnement :} Économétrie, DSGE, Python, R, Excel, Julia, GAMS.
\end{itemize}
\divider
\cvevent{Enseignant vacataire -- Pauvreté \& Conditions de Vie des Ménages + Enquête Sociologique}{Institut Supérieur des Métiers de la Statistique (ISS)}{Février 2026 -- Juin 2026}{Nouakchott, Mauritanie}
\begin{itemize}
  \item Enseignement du module Pauvreté et Conditions de Vie des Ménages (édition améliorée).
  \item Conception et enseignement d'un nouveau module : Enquête Sociologique (méthodologie d'enquête, échantillonnage, conception de questionnaires, analyse des données).
\end{itemize}
\divider
\cvevent{Stage de Fin d'Étude}{Caisse des Dépôts et de Développement (CDD)}{Février -- Juin 2023}{Nouakchott, Mauritanie}
\begin{itemize}
  \item Analyse des données bancaires issues de la CDD pour évaluer l'impact du système de garantie de crédit sur les défauts de paiement en Mauritanie.
  \item Utilisation de modèles économétriques et d'approches modernes de machine learning pour prédire les risques de défaut.
\end{itemize}
\divider
\cvevent{Stage Ouvrier}{Agence Nationale de la Statistique et de l'Analyse Démographique et Économique (ANSADE)}{Avril -- Juin 2022}{Nouakchott, Mauritanie}
\begin{itemize}
  \item Élaboration d'une méthodologie pour la cartographie de données en vue de la conception des comptes sectoriels et analyse sous ERETES.
\end{itemize}

\cvsection{Projets}
\cvevent{Modeling the Drivers of Mauritania's GDP (ARDL)}{Modélisation économétrique de la croissance}{}{}
\begin{itemize}
  \item Modélisation économétrique de la dynamique de croissance de la Mauritanie à l'aide d'un modèle ARDL (AutoRegressive Distributed Lag) sous R. \href{https://github.com/MoustaphaCheikhJidou/ARDL-Growth-Mauritania}{GitHub}.
\end{itemize}
\divider
\cvevent{Poverty \& Inequality Forecasting in Mauritania}{Prévisions d'indicateurs de pauvreté}{}{}
\begin{itemize}
  \item Prévisions de quatre indicateurs clés de pauvreté de 1996 à 2030 à partir de données d'enquête et de techniques d'interpolation pour appuyer la planification des politiques publiques. \href{https://github.com/MoustaphaCheikhJidou/Poverty-Gini-Forecast-Mauritania}{GitHub}.
\end{itemize}

\switchcolumn \cvsection{Formation}
\cvevent{Faculté des Sciences Économiques et de Gestion}{FSEG - Mauritanie}{Octobre 2023 -- Juin 2025}{} Master en Économétrie et Statistiques Appliquées.
\divider
\cvevent{École Supérieure Polytechnique de Nouakchott}{ESP - Mauritanie}{Octobre 2023 -- Juin 2026}{} Ingénieur d'État en Statistique et en Ingénierie des Données.

\cvsection{Compétences}
{\footnotesize R \textbullet{} GAMS \textbullet{} Eviews \textbullet{} Julia \textbullet{} Excel \textbullet{} Modélisation DSGE \textbullet{} Analyse de politique monétaire \textbullet{} Comptabilité nationale}

\cvsection{Langues} \cvskill{Arabe}{4}{Langue natale}\\ \divider \cvskill{Français}{3}{DALF C1}\\ \divider \cvskill{Anglais}{2}{CEFR B1}

\cvsection{Vie Associative}
{\footnotesize Président de l'Alumni de l'ISS}

\cvsection{Certifications}
{\footnotesize \textbf{Google} Data Analytics \textbullet{} \textbf{DataCamp} Intermediate R}

\end{paracol}
\end{document}
